\chapter{Membaca Kode Kehidupan}
\label{ch:genomik}

Pada musim panas 2025 saya membuka pintu ke ilmu hayati lewat beberapa kuliah sekaligus: pengantar
genetika, analisis genom dan transkriptomik, bioinformatika struktural, dan AI dalam ilmu hayati.
Kuliah genom dimulai dari dasar sekali, dari apa itu kehidupan dan bagaimana DNA disusun, lalu
bergerak ke teknologi microarray, pengurutan DNA generasi demi generasi, perakitan genom, pemetaan
bacaan, pemanggilan varian, dan analisis RNA-seq. Bagi pembaca yang datang dari teknik atau ilmu
komputer, seperti saya, bab ini adalah perkenalan dengan bahasa yang dipakai bab-bab berikutnya.

Biologi molekuler modern adalah ilmu data. Satu genom manusia berisi sekitar tiga miliar huruf,
satu percobaan pengurutan menghasilkan ratusan juta potongan bacaan, dan satu pengukuran ekspresi
gen mencatat aktivitas puluhan ribu gen sekaligus. Hampir setiap langkah dari sampel ke kesimpulan
melibatkan algoritme, statistik, dan semakin sering machine learning.

\section{Dogma sentral}

Informasi genetik tersimpan dalam DNA, rantai ganda yang tersusun dari empat basa: adenin, sitosin,
guanin, dan timin. Basa pada satu untai berpasangan dengan basa pada untai lain, A dengan T dan C
dengan G, sehingga setiap untai menyimpan salinan lengkap informasi untai pasangannya. Struktur
heliks ganda ini diusulkan Watson dan Crick pada 1953, dengan bantuan penting dari foto difraksi
sinar-X yang dibuat Rosalind Franklin, sebuah kisah yang juga disinggung di kuliah.

Dogma sentral biologi molekuler menggambarkan aliran informasi. DNA disalin menjadi RNA dalam proses
transkripsi, lalu RNA pembawa pesan diterjemahkan menjadi protein di ribosom. Setiap tiga basa, satu
kodon, menyandikan satu asam amino. Ada $4^3=64$ kodon untuk dua puluh asam amino dan sinyal berhenti,
sehingga kodenya berlebih: beberapa kodon berbeda menyandikan asam amino yang sama. Pada sel eukariot,
bagian-bagian RNA dipotong dan disambung ulang sebelum diterjemahkan, dan penyambungan alternatif
memungkinkan satu gen menghasilkan beberapa protein berbeda. Protein lalu melipat menjadi bentuk tiga
dimensi yang menentukan fungsinya, topik \cref{ch:protein}.

Bayangkan perpustakaan besar yang buku-bukunya tidak boleh dibawa keluar. Untuk membangun sesuatu,
seseorang memfotokopi satu halaman resep (transkripsi), membawa fotokopinya ke dapur (ribosom), lalu
memasak sesuai resep itu (translasi). Sel yang berbeda di tubuh kita membawa perpustakaan yang sama,
tetapi memfotokopi halaman yang berbeda. Itulah sebabnya sel otot dan sel saraf begitu berbeda
meski DNA-nya identik, dan itulah sebabnya mengukur RNA, yaitu fotokopi yang sedang beredar, memberi
tahu apa yang sedang dikerjakan sebuah sel.

\section{Variasi genetik}

Dua manusia berbeda di sekitar satu dari setiap seribu basa. Variasi yang paling umum adalah
\istilah{single nucleotide polymorphism} (SNP), satu posisi tempat sebagian orang punya satu basa dan
sebagian lain punya basa berbeda. Genetika populasi mempelajari bagaimana frekuensi varian berubah.
Hukum Hardy--Weinberg menyatakan bahwa dalam populasi besar yang kawin acak tanpa seleksi, kalau
frekuensi dua alel adalah $p$ dan $q=1-p$, maka frekuensi genotipenya $p^2$, $2pq$, dan $q^2$. Dengan
$p=0{,}7$, sekitar 49 persen individu homozigot untuk alel pertama, 42 persen heterozigot, dan 9 persen
homozigot untuk alel kedua. Penyimpangan dari proporsi ini menjadi petunjuk adanya seleksi, migrasi,
atau kesalahan pengukuran.

Varian yang letaknya berdekatan di kromosom cenderung diwariskan bersama, sebuah keadaan yang disebut
\istilah{linkage disequilibrium}. Studi asosiasi seluruh genom (GWAS) memanfaatkan hal ini: mereka mengukur
ratusan ribu SNP pada ribuan orang dan mencari varian yang frekuensinya berbeda antara penderita dan
bukan penderita sebuah penyakit. Karena jumlah uji statistiknya banyak sekali, ambang signifikansinya
harus ketat sekali.

\section{Mengukur ekspresi gen}

Microarray DNA adalah kaca kecil yang ditempeli jutaan potongan DNA pendek yang dikenal, yang disebut
probe. RNA dari sampel diubah menjadi DNA berlabel fluoresen, lalu dibiarkan menempel pada probe yang
cocok. Semakin terang sebuah titik, semakin banyak RNA dari gen itu. Jalan dari gambar ke angka
ternyata panjang: koreksi latar belakang, normalisasi antar chip agar distribusinya sebanding,
misalnya dengan normalisasi kuantil, lalu meringkas banyak probe dari satu gen menjadi satu nilai
ekspresi. Kuliah membahas metode ringkasan seperti RMA, dan FARMS, metode berbasis analisis faktor
(\cref{ch:unsup}) dari kelompok Sepp Hochreiter, yang memperlakukan sinyal bersama antar probe sebagai faktor
laten dan noise pribadi setiap probe sebagai noise analisis faktor.

\section{Pengurutan DNA}

Pengurutan Sanger, generasi pertama, membaca satu potongan DNA sepanjang beberapa ratus basa dengan
akurasi tinggi, dan dipakai dalam Proyek Genom Manusia. Generasi kedua, yang dipelopori platform
Illumina, membaca ratusan juta potongan pendek secara paralel, sering berpasangan dari kedua ujung
sebuah fragmen. Generasi ketiga dan keempat membaca satu molekul DNA secara langsung, dengan bacaan
yang jauh lebih panjang, dari puluhan ribu sampai jutaan basa, meski dulu dengan galat yang lebih tinggi.
Pada pengurutan nanopore, DNA ditarik melewati pori kecil dan setiap kombinasi basa mengubah arus
listrik dengan cara yang khas. Mengubah sinyal arus menjadi urutan basa, \istilah{base calling}, kini
dikerjakan jaringan saraf rekuren dan konvolusi yang dilatih dengan loss connectionist temporal
classification, karena panjang keluaran tidak diketahui sebelumnya.

Setiap basa yang dibaca diberi skor kualitas PHRED, $Q=-10\log_{10}P_{\mathrm{galat}}$. Skor 30 berarti
peluang salah satu per seribu. Kedalaman cakupan, $C=NL/G$, menyatakan berapa kali rata-rata setiap
posisi genom dibaca, dengan $N$ jumlah bacaan, $L$ panjang bacaan, dan $G$ panjang genom. Untuk genom
manusia sepanjang tiga miliar basa, cakupan tiga puluh kali dengan bacaan 150 basa membutuhkan sekitar
600 juta bacaan.

\section{Dari bacaan ke genom}

Kalau genom acuan sudah ada, setiap bacaan dipetakan ke posisinya di acuan. Mencocokkan ratusan juta
bacaan pendek ke tiga miliar huruf dengan toleransi terhadap galat membutuhkan indeks yang cerdas.
Transformasi Burrows--Wheeler, yang juga dipakai dalam kompresi data, memungkinkan pencarian pola yang
amat cepat dengan memori kecil \citep{li2009}. Strategi umumnya adalah \istilah{seed and extend}: cari
potongan pendek yang cocok persis, lalu perluas dengan penjajaran yang toleran terhadap galat. Untuk
bacaan panjang, potongan-potongan yang cocok dirangkai dulu menjadi rantai yang konsisten sebelum
penjajaran rinci.

Setelah bacaan dipetakan, \istilah{variant calling} memutuskan di posisi mana sampel berbeda dari acuan.
Tantangannya membedakan varian sungguhan dari galat pengurutan dan galat pemetaan. DeepVariant
\citep{poplin2018} mengubah tumpukan bacaan di sekitar sebuah posisi menjadi gambar dan memakai jaringan
konvolusi untuk mengklasifikasikan genotipenya, sebuah contoh menarik dari visi komputer yang dipinjam
oleh genomik.

Kalau tidak ada acuan, genom harus dirakit dari nol, seperti menyusun teka-teki jigsaw raksasa dari
potongan-potongan yang saling tumpang tindih. Pendekatan \istilah{overlap--layout--consensus} membangun graf
yang simpulnya bacaan dan sisinya tumpang tindih, lalu mencari jalur yang melewati setiap simpul tepat
sekali, sebuah jalur Hamilton yang sulit dicari secara umum. Graf de Bruijn membalik sudut pandangnya:
bacaan dipotong menjadi semua kata sepanjang $k$, dan rakitan menjadi jalur yang melewati setiap sisi
tepat sekali, sebuah jalur Euler yang bisa dicari dengan efisien. Bagian yang berulang di genom adalah
musuh utama keduanya, karena potongan dari dua tempat berbeda tampak identik. Bacaan panjang membantu
justru karena mereka bisa melompati pengulangan. Kuliah juga menunjukkan bahwa jaringan graf
(\cref{ch:gdl}) mulai dipakai untuk menyederhanakan graf rakitan.

\section{RNA-seq}

Pengurutan juga bisa mengukur ekspresi gen. RNA diubah menjadi DNA, diurutkan, lalu jumlah bacaan
yang jatuh pada setiap gen dihitung. Hitungan ini harus dinormalisasi karena sampel yang berbeda
diurutkan dengan kedalaman berbeda. Untuk mencari gen yang ekspresinya berubah antara dua kondisi,
DESeq2 \citep{love2014} memodelkan hitungan dengan distribusi binomial negatif, yang mengizinkan varians
lebih besar dari rata-rata seperti yang selalu terlihat pada data biologis. Karena setiap gen hanya
punya sedikit ulangan, perkiraan dispersi dan perubahan lipat dari banyak gen ``dipinjamkan'' satu sama
lain dan ditarik ke arah tren bersama, sebuah bentuk prior Bayesian empiris. Dengan puluhan ribu gen
yang diuji sekaligus, nilai p dikoreksi untuk mengendalikan proporsi temuan palsu.

\section{Di mana machine learning masuk}

Genomik adalah salah satu medan tempat machine learning paling banyak dipakai di ilmu hayati. Jaringan
saraf memanggil basa dari sinyal nanopore, memanggil varian dari tumpukan bacaan, dan memprediksi dari
urutan DNA saja di mana protein akan menempel atau seberapa kuat sebuah gen diekspresikan. Model bahasa
yang dilatih pada urutan DNA dan protein meniru resep model bahasa manusia di \cref{ch:lstm}.
Metode tanpa label seperti biclustering dan analisis faktor (\cref{ch:unsup}) menemukan modul gen yang
bekerja bersama dalam data ekspresi. Karakter datanya menuntut kehati-hatian khusus: jumlah fitur
puluhan ribu dengan jumlah sampel puluhan, batch percobaan yang meninggalkan jejak sistematis, dan
korelasi antar sampel dari pasien atau keluarga yang sama, semua sumber kebocoran yang dibahas di
\cref{ch:data}.

\sumberbab{Bab ini mengikuti kuliah Genome Analysis and Transcriptomics (dasar biologi molekuler,
genetika populasi, microarray dan pemrosesannya, pengurutan DNA, pemetaan bacaan, pemanggilan varian,
perakitan genom, RNA-seq) dan Introduction to Genetics. Untuk analisis urutan secara probabilistik,
\citet{durbin1998} tetap menjadi rujukan klasik.}
